\documentclass[a4paper,11pt]{article}
\usepackage[utf8]{inputenc}
\usepackage[T1]{fontenc}
\usepackage[margin=1in]{geometry}

% Coordinated academic typography
\usepackage{newtxtext}
\usepackage{newtxmath}
\usepackage{microtype}

\usepackage{graphicx}
\usepackage{booktabs}
\usepackage{float}
\usepackage{url}
\usepackage{amsmath}

% Slightly cleaner paragraph spacing for a short technical report
\setlength{\parindent}{0pt}
\setlength{\parskip}{0.55em}

\title{Short Benchmark Report}
\author{Bin Zhou, Zekai Lin, Yu Zhang}
\date{13 September 2026}

\begin{document}

\maketitle

\section{Storage Strategies Evaluated}

The benchmark evaluates four main Delta Lake storage strategies for the Taxi Trips dataset:

\begin{enumerate}
    \item Unpartitioned
    \item Partitioned by pickup\_borough
    \item Partitioned by pickup\_month
    \item Partitioned by pickup\_date
\end{enumerate}

Three additional random file-count control strategies were evaluated to investigate the role of physical file layout:

\begin{itemize}
    \item random\_20\_files, matching the 20 Parquet files generated by pickup\_month;
    \item random\_128\_files, matching the 128 Parquet files generated by pickup\_borough;
    \item random\_608\_files, matching the 608 Parquet files generated by pickup\_date.
\end{itemize}

The controls assign records to deterministic pseudo-random buckets using a hash of the trip identifier and a fixed seed. They are experimental comparison layouts, rather than proposed analytical partitioning schemes. Matching file counts helps interpret the results, although it does not hold file contents, compression, and task distribution constant.

\subsection*{Dataset and Methodology}

All strategies use the same projection of the integrated Delta table, containing 9,551,387 cleaned taxi trips from January--March 2024. The projection retains the taxi attributes required by the benchmark, pickup borough, and source-schema lineage. The lineage snapshot records version 1 for all contributing source datasets. Weather and air-quality measurements are excluded from the benchmark projection.

The results reported here are from the benchmark run completed on 13 September 2026. The Delta commit metadata identifies Apache Spark 3.5.9 and Delta Lake 3.3.3. Execution was local on Windows. The CSV does not record the hardware specification or effective runtime overrides, which limits exact reproduction of the absolute timings on another machine.

For each strategy, the table was written three times, removing the previous benchmark table directory before each write. The median write time is reported, and the final written table is retained for querying. Write latency includes reading the projected integrated input and executing the transformations required by the write. It does not measure the complete source-file ingestion, cleaning, and integration pipeline. Random layouts additionally perform explicit repartitioning, so their write times include that shuffle cost.

Each query was executed once as a warm-up and then measured five times, using \texttt{collect()} to trigger execution. The median query latency is reported. Spark's table cache was cleared before measured query runs, but the operating-system file cache was not explicitly cleared. These are therefore repeated, warmed-access measurements rather than controlled cold-disk measurements. Strategies were executed in a fixed order; background activity and cache state may contribute to small timing differences.

The following three queries required by Task 6 were evaluated:

\begin{itemize}
    \item number of taxi trips per pickup borough;
    \item average trip duration in seconds per pickup date;
    \item average fare amount per pickup borough.
\end{itemize}

All three queries are full-table aggregations without partition-filter predicates. They evaluate storage layouts for this workload and do not measure the benefits of partition pruning for selective queries.

\section{Benchmark Results}

\begin{table}[H]
\centering
\small
\resizebox{\textwidth}{!}{
\begin{tabular}{lrrrrrr}
\toprule
Strategy & Write median (s) & Storage (MiB) & Parquet files & Trips/borough (s) & Duration/day (s) & Fare/borough (s) \\
\midrule
Unpartitioned & 7.8079 & 788.33 & 16 & 0.3649 & 0.3089 & 0.3969 \\
\texttt{pickup\_month} & 7.6955 & 788.11 & 20 & 0.3230 & \textbf{0.2619} & \textbf{0.3094} \\
Random 20 files & 14.3658 & 800.19 & 20 & \textbf{0.3132} & 0.2814 & 0.3373 \\
\texttt{pickup\_borough} & \textbf{7.5985} & 785.98 & 128 & 0.3160 & 0.3323 & 0.3178 \\
Random 128 files & 11.4217 & 809.40 & 128 & 0.3681 & 0.3444 & 0.3866 \\
\texttt{pickup\_date} & 9.1207 & 798.14 & 608 & 0.5874 & 0.5309 & 0.6216 \\
Random 608 files & 18.8582 & 830.72 & 608 & 0.5702 & 0.5668 & 0.6156 \\
\bottomrule
\end{tabular}
}
\caption{Latest benchmark results. Bold values indicate the lowest median for each timing metric across all seven strategies.}
\label{tab:benchmark_results}
\end{table}

Storage is the total table-directory size divided by $1024^2$, including data files and auxiliary files such as Delta logs and local checksum files. The unit is therefore MiB. The file-count column counts Parquet data files only. Removing each benchmark directory before rewriting prevents retained historical Delta data files from inflating this measurement.

\subsection*{Result Validation}

All seven stored tables contain 9,551,387 rows. A separate verification reran the three required queries on each layout. Grouping keys, trip counts, and average trip durations matched across all layouts. Differences in average fare were below $10^{-12}$ when compared with the borough-partitioned layout, consistent with floating-point accumulation in different execution orders. These differences are below an absolute comparison tolerance of $10^{-9}$ and do not change the analytical interpretation.

\section{Performance Discussion}

Among the four practical storage strategies, pickup\_borough has the lowest observed median write time, 7.5985 seconds, and the fastest trips-per-borough query, 0.3160 seconds. Pickup\_month has the fastest duration-per-day and fare-per-borough queries, at 0.2619 and 0.3094 seconds respectively. It also generates only 20 Parquet files and is faster than the unpartitioned layout for all three queries in this run. Thus, monthly partitioning is a reasonable overall compromise, but it is not the best strategy for every measured metric.

The write-time difference between the borough and monthly layouts is approximately 0.097 seconds, or 1.3\%. The trips-per-borough difference is only 0.007 seconds. These rankings describe the observed medians; the small differences should not be interpreted as robust performance advantages without additional independent repetitions.

Partitioning by pickup\_date is consistently slower for the measured queries. It produces 608 Parquet files, with query latencies of approximately 0.53--0.62 seconds. These are about 1.82--2.03 times the corresponding monthly-layout latencies. This pattern is consistent with additional file-opening, scheduling, and metadata overhead from a more fragmented layout, although these individual costs were not measured separately.

The random controls support a similar interpretation: random\_20\_files completes the queries in approximately 0.28--0.34 seconds, whereas random\_608\_files requires approximately 0.57--0.62 seconds. However, random layouts do not consistently outperform semantic layouts with the same file count. Random 20 is slightly faster than pickup\_month for trip counts but slower for both averages. Random 128 is slower than pickup\_borough for all three queries. Random 608 is slightly faster than pickup\_date for trip counts and average fare, but slower for average duration.

Consequently, file count is useful for understanding the observed pattern, but it is not the only relevant property. Row distribution, file size, compression, and execution variability may also affect performance. Without task-level measurements, the results do not establish that reduced skew explains any particular timing difference.

Random-layout write times range from approximately 11.42 to 18.86 seconds and include the extra shuffle used to control the layout. They represent the cost of constructing those experimental layouts, rather than an isolated comparison of partition-column choices.

Storage sizes for the four main strategies are relatively close, ranging from 785.98 to 798.14 MiB. In this experiment, query latency varies more substantially than storage size. The results therefore highlight the importance of physical layout for repeated full-table aggregation, rather than a large storage-capacity advantage for one scheme.

\section{Conclusion}

For this dataset and the three full-table analytical queries, daily partitioning produces many small files and consistently higher query latency than monthly partitioning. The random file-count controls provide supporting evidence for a file-layout effect, but do not isolate file fragmentation as the sole cause.

Pickup\_month is a reasonable default among the tested practical layouts: it combines a low file count with the fastest two average queries and a write time close to the fastest observed strategy. Pickup\_borough remains competitive and is slightly faster than pickup\_month for the trips-per-borough query in this run. The choice should therefore reflect the expected query mix rather than a claim that one layout is universally optimal.

These conclusions apply to the measured local workload. Since the queries contain no partition filters, the experiment cannot determine the benefit of date-, month-, or borough-based pruning for selective queries. Such workloads, additional data volumes, and independent repeated runs would be needed before generalizing the storage recommendation.

\end{document}
